\subsection*{Does $P$(post-exilic) mean what it says?}
\label{sec:postexilic_cal}

The errors in Table~\ref{tab:leakage} are large: the likelihood-only estimate
misses by 156~yr on average and by up to 334~yr on individual units, and
there is no detectable improvement with text length
($\rho = -0.23$ between $|$error$|$ and word count, $p = 0.27$).  A reader
who stops there will conclude that no claim resting on this model can be
trusted, and the conclusion would be correct for any claim about a
\emph{date}.

The substantive claim of this paper is not about a date.  It is that a unit's
predictive mass lies on one side of a single boundary, and that is a much
coarser question.  Whether it is answerable is an empirical matter, so we
test it directly rather than argue for it.

\subsection*{Protocol}

For each of the 25 dated units the model is fitted on the other 24, and the
conformal residual set is computed by a second, inner leave-one-out
\emph{within} those 24.  The held-out unit therefore contributes neither to
its point estimate nor to the residuals that give that estimate its spread.
The predictive distribution is the point estimate displaced by that residual
set, and $P_{\text{post}}$ is the fraction of it later than 586~BCE.  Each
unit's true side of the boundary is then compared with the probability the
model assigned before seeing it.

\subsection*{Result}

The probabilities are calibrated (Fig~\ref{fig:calibration}).  In the
highest bin the generative family predicts a mean 0.95 and 100\% of those
units are post-exilic; ridge predicts 0.92 against an observed 0.89.  Of the
units on which either family is confident --- $P_{\text{post}} \geq 0.8$ or
$\leq 0.2$ --- the generative family is correct on 9 of 9 and ridge on 8 of
9.  Brier skill scores against the corpus base rate are $+0.200$ and
$+0.187$; under a permutation null in which the true labels are shuffled,
$p = 0.005$ and $p = 0.014$.

Restricting to units more than 200~yr from the boundary --- the regime the
P source occupies, at 242--322~yr --- the generative family is correct on 5
of 5 confident calls and ridge on 4 of 4.

The two findings are the same finding.  Because the typical error is of
order 150~yr, the model reaches $P_{\text{post}} \geq 0.9$ only when a unit
sits well over 200~yr from the boundary.  The large errors of
Table~\ref{tab:leakage} are not in tension with the probability of
Table~\ref{tab:targets}: they are what produces it.  A 334~yr error on
Haggai and a 0.92 on the P source are consequences of the same residual
distribution, and the second already accounts for the first.

\subsection*{Where this test cannot reach}

Three limits should be stated plainly.

First, one confident call is wrong.  Ridge assigns $P_{\text{post}} = 1.00$
to Nahum, which is pre-exilic on a secure terminus (before the fall of
Nineveh in 612~BCE).  That is one error in 25, it is not shared by the
generative family, and it is the kind of failure a corpus this size cannot
rule out.

Second, units close to the boundary are decided by arithmetic rather than by
evidence.  Obadiah is dated 585~BCE, one year into the post-exilic side; both
families place it post-exilic and are scored correct, but no method
distinguishing 585 from 587 on linguistic grounds is doing real work.
Excluding units within 25~yr of the boundary leaves the generative family at
7 of 7 and ridge at 3 of 4.

Third, and most consequentially, the test is asymmetric and cannot be made
otherwise.  The earliest dated unit in the corpus is Amos at 760~BCE, only
174~yr before the boundary, while the P source estimate lies 242--322~yr
\emph{after} it.  Confident post-exilic calls can therefore be validated at
the distance the claim requires; confident pre-exilic calls at that distance
cannot be, because no Hebrew text of sufficient length exists that early.
This is not a limitation of corpus selection that further work could
remedy.  Morphologically tagged Hebrew of the kind this method requires
exists for the Masoretic text and essentially nowhere else, and the
non-biblical corpora that might one day be tagged --- Qumran, Ben Sira ---
are themselves late, so they would deepen the asymmetry rather than correct
it.  Any linguistic dating of Biblical Hebrew inherits this constraint.  It
bears directly on how the present results should be read: the framework can
support a claim that a text is later than a boundary far better than it can
support a claim that a text is earlier than one.
